\chapter{数学的準備}
\label{ch:seminar-00}

%% ============================================================
%%  0. 置換と置換行列
%% ============================================================
\section{置換と置換行列}
\label{sec:sem-permutation}

\begin{definition}{置換}{sem-permutation}
  $\range{n} \defeq \{1, \ldots, n\}$ 上の\textbf{置換}（permutation）とは，
  $\range{n}$ から $\range{n}$ への全単射をいう．
  $\range{n}$ 上のすべての置換の集合を $\Perm(n)$ と記す．
  $|\Perm(n)| = n!$ である．
\end{definition}

\begin{definition}{置換行列}{sem-permutation-matrix}
  置換 $\sigma \in \Perm(n)$ に対して，
  $n \times n$ 行列 $\mathbf{P}_\sigma$ を
  \[
    (\mathbf{P}_\sigma)_{i,j} =
    \begin{cases}
      1/n & j = \sigma(i), \\
      0 & j \neq \sigma(i)
    \end{cases}
  \]
  で定義し，これを\textbf{置換行列}（permutation matrix）という．
  各行・各列にちょうど1つの非零成分 $1/n$ を持ち，
  $\mathbf{P}_\sigma \ones_n = \ones_n / n$，
  $\mathbf{P}_\sigma^\top \ones_n = \ones_n / n$
  が成り立つ．
\end{definition}

\begin{example}{置換行列}{sem-permutation-matrix-example}
  $n = 3$，置換 $\sigma = (2, 3, 1)$
  （すなわち $\sigma(1) = 2,\; \sigma(2) = 3,\; \sigma(3) = 1$）のとき，
  置換行列は
  \[
    \mathbf{P}_\sigma =
    \frac{1}{3}
    \begin{pmatrix}
      0 & 1 & 0 \\
      0 & 0 & 1 \\
      1 & 0 & 0
    \end{pmatrix}.
  \]
  行和・列和を確認すると，
  $\mathbf{P}_\sigma \ones_3
  = \frac{1}{3} \begin{psmallmatrix} 1 \\ 1 \\ 1 \end{psmallmatrix}
  = \ones_3 / 3$．
\end{example}

%% ============================================================
%%  1. フロベニウス内積
%% ============================================================
\section{フロベニウス内積}
\label{sec:sem-frobenius-inner}

\begin{definition}{フロベニウス内積}{sem-frobenius-inner}
  同じサイズの2つの行列 $\mathbf{A}, \mathbf{B} \in \R^{n \times m}$ に対し，
  \textbf{フロベニウス内積}（Frobenius inner product）を
  \[
    \inner{\mathbf{A}}{\mathbf{B}}
    \defeq \tr(\mathbf{A}^\top \mathbf{B})
    = \sum_{i=1}^{n} \sum_{j=1}^{m} A_{i,j} B_{i,j}
  \]
  で定める．
\end{definition}

\begin{example}{フロベニウス内積の計算}{sem-frobenius-inner-example}
  \[
    \mathbf{A} =
    \begin{pmatrix} 1 & 2 \\ 3 & 4 \end{pmatrix}, \quad
    \mathbf{B} =
    \begin{pmatrix} 5 & 6 \\ 7 & 8 \end{pmatrix}
  \]
  に対し，成分ごとの積の総和から
  \[
    \inner{\mathbf{A}}{\mathbf{B}}
    = 1 \cdot 5 + 2 \cdot 6 + 3 \cdot 7 + 4 \cdot 8
    = 5 + 12 + 21 + 32
    = 70.
  \]
  トレースによる計算でも確かめられる：
  \[
    \mathbf{A}^\top \mathbf{B}
    = \begin{pmatrix} 1 & 3 \\ 2 & 4 \end{pmatrix}
      \begin{pmatrix} 5 & 6 \\ 7 & 8 \end{pmatrix}
    = \begin{pmatrix} 26 & 30 \\ 38 & 44 \end{pmatrix}, \quad
    \tr(\mathbf{A}^\top \mathbf{B}) = 26 + 44 = 70.
  \]
\end{example}

%% ============================================================
%%  2. 距離空間
%% ============================================================
\section{距離空間}
\label{sec:sem-metric-space}

\begin{definition}{距離空間}{sem-metric-space}
  集合 $\Omega$ と関数 $d \colon \Omega \times \Omega \to [0, \infty)$ の組
  $(\Omega, d)$ が\textbf{距離空間}（metric space）であるとは，
  次の3条件がすべての $x, y, z \in \Omega$ に対して成り立つことをいう：
  \begin{enumerate}
    \item[(i)] 正定値性：$d(x, y) = 0 \iff x = y$，
    \item[(ii)] 対称性：$d(x, y) = d(y, x)$，
    \item[(iii)] 三角不等式：$d(x, z) \leq d(x, y) + d(y, z)$．
  \end{enumerate}
  関数 $d$ を $\Omega$ 上の\textbf{距離関数}という．
\end{definition}

\begin{remark}{開球と開集合}{sem-open-ball}
  距離空間 $(\Omega, d)$ 上の点 $x \in \Omega$ と半径 $r > 0$ に対して，
  \textbf{開球} $B_r(x)$ を
  \[
    B_r(x) \defeq \{y \in \Omega \mid d(x, y) < r\}
  \]
  で定める．
  $\Omega$ の部分集合 $U \subseteq \Omega$ が\textbf{開集合}であるとは，
  任意の $x \in U$ に対してある $r > 0$ が存在して
  $B_r(x) \subseteq U$ となることをいう．
  開集合の補集合を\textbf{閉集合}という．
\end{remark}

%% ============================================================
%%  3. 可測空間と確率測度
%% ============================================================
\section{可測空間と確率測度}
\label{sec:sem-measurable}

\begin{definition}{$\sigma$-代数}{sem-sigma-algebra}
  集合 $\Omega$ の部分集合族 $\mathcal{F} \subseteq 2^{\Omega}$ が
  \textbf{$\sigma$-代数}（$\sigma$-algebra）であるとは，
  以下の3条件を満たすことをいう：
  \begin{enumerate}
    \item[(i)] $\Omega \in \mathcal{F}$，
    \item[(ii)] $A \in \mathcal{F} \Longrightarrow A^c \in \mathcal{F}$（補集合で閉じる），
    \item[(iii)] $A_1, A_2, \ldots \in \mathcal{F} \Longrightarrow
      \bigcup_{k=1}^{\infty} A_k \in \mathcal{F}$（可算合併で閉じる）．
  \end{enumerate}
\end{definition}

\begin{definition}{可測空間}{sem-measurable-space}
  集合 $\Omega$ と $\Omega$ 上の $\sigma$-代数 $\mathcal{F}$ の組
  $(\Omega, \mathcal{F})$ を\textbf{可測空間}（measurable space）という．
  $\mathcal{F}$ の元を\textbf{可測集合}という．
\end{definition}

\begin{definition}{可測写像}{sem-measurable-map}
  可測空間 $(\Omega_1, \mathcal{F}_1)$ から $(\Omega_2, \mathcal{F}_2)$ への写像
  $T \colon \Omega_1 \to \Omega_2$ が\textbf{可測}であるとは，
  $T^{-1}(A) \in \mathcal{F}_1$ が
  すべての $A \in \mathcal{F}_2$ に対して成り立つことをいう．
\end{definition}

\begin{definition}{ボレル $\sigma$-代数}{sem-borel}
  距離空間 $(\Omega, d)$ において，
  $d$ から定まる距離位相の開集合すべてを含む最小の $\sigma$-代数を
  \textbf{ボレル $\sigma$-代数}といい，$\Bb(\Omega)$ と記す．
\end{definition}

\begin{definition}{測度}{sem-measure}
  可測空間 $(\Omega, \mathcal{F})$ 上の\textbf{測度}（measure）とは，
  関数 $\mu \colon \mathcal{F} \to [0, \infty]$ であって，
  \begin{enumerate}
    \item[(i)] $\mu(\emptyset) = 0$，
    \item[(ii)] 互いに素な可測集合 $A_1, A_2, \ldots \in \mathcal{F}$ に対して
      $\mu\bigl(\bigcup_{k=1}^{\infty} A_k\bigr) = \sum_{k=1}^{\infty} \mu(A_k)$
      （可算加法性）
  \end{enumerate}
  を満たすものをいう．
\end{definition}

\begin{definition}{確率測度}{sem-probability-measure}
  可測空間 $(\Omega, \mathcal{F})$ 上の測度 $\mu$ が
  $\mu(\Omega) = 1$ を満たすとき，$\mu$ を
  \textbf{確率測度}（probability measure）という．
  $\Omega$ 上の確率測度の全体を $\Mm_+^1(\Omega)$ と記す．
\end{definition}

%% ============================================================
%%  4. 本稿で用いる可測空間
%% ============================================================
\section{本稿で用いる可測空間 $(\X, \Bb(\X))$}
\label{sec:sem-borel-measurable-space}

本節では，前節までの一般論を本稿の舞台となる距離空間 $\X, \Y$ に適用し，
以降で常用する可測空間の取り決め・用語・性質をまとめる．

\begin{remark}{本稿における可測構造の取り決め}{sem-borel-convention}
  可測空間は集合と $\sigma$-代数の組であるから，距離空間 $\X$ を可測空間として
  扱うには $\X$ 上に $\sigma$-代数を指定する必要がある．
  本稿では以降，距離空間 $\X$ に付随する $\sigma$-代数として
  常にボレル $\sigma$-代数 $\Bb(\X)$ を採用し，$(\X, \Bb(\X))$ を既定の可測空間とする．
  $\Y$ についても同様である．

  この取り決めのもとで用語を次のように約束する：
  \begin{itemize}
    \item 「$\X$ の\textbf{可測集合}」とは $\Bb(\X)$ の元（\textbf{ボレル集合}）を指す．
    \item 「$\X$ 上の確率測度」とは $(\X, \Bb(\X))$ 上の確率測度
      $\mu \in \Mm_+^1(\X)$ を指す．
    \item 距離空間 $\X, \Y$ の間の写像 $T \colon \X \to \Y$ が
      \textbf{ボレル可測}であるとは，$(\X, \Bb(\X))$ から $(\Y, \Bb(\Y))$ への
      可測写像であることをいう．
  \end{itemize}
\end{remark}

\begin{proposition}{$(\X, \Bb(\X))$ の基本性質}{sem-borel-properties}
  距離空間 $\X$ とボレル $\sigma$-代数 $\Bb(\X)$ に対し，次が成り立つ：
  \begin{enumerate}
    \item[(i)] 補集合で閉じる（$\sigma$-代数の公理，定義~\ref{def:sem-sigma-algebra} (ii)）：
      \[
        A \in \Bb(\X) \;\Longrightarrow\; A^c \in \Bb(\X).
      \]
    \item[(ii)] 可算合併で閉じる（$\sigma$-代数の公理，定義~\ref{def:sem-sigma-algebra} (iii)）：
      任意の列 $A_1, A_2, \ldots \in \Bb(\X)$ に対して
      \[
        \bigcup_{k=1}^{\infty} A_k \in \Bb(\X).
      \]
    \item[(iii)] 可算共通部分で閉じる（(i), (ii) と De Morgan の法則
      $\bigcap_{k} A_k = \bigl(\bigcup_{k} A_k^c\bigr)^c$ から従う）：
      \[
        A_1, A_2, \ldots \in \Bb(\X)
        \;\Longrightarrow\;
        \bigcap_{k=1}^{\infty} A_k \in \Bb(\X).
      \]
    \item[(iv)] 開集合はすべてボレル集合である
      （ボレル $\sigma$-代数の定義~\ref{def:sem-borel}）：
      \[
        U \subseteq \X \text{ が開集合}
        \;\Longrightarrow\; U \in \Bb(\X).
      \]
    \item[(v)] 閉集合はすべてボレル集合である（(i), (iv) と $F = (F^c)^c$ から従う）：
      \[
        F \subseteq \X \text{ が閉集合}
        \;\Longrightarrow\; F \in \Bb(\X).
      \]
    \item[(vi)] ボレル可測写像の引き戻しはボレル集合である
      （可測写像の定義~\ref{def:sem-measurable-map} を
      $(\X, \Bb(\X)) \to (\Y, \Bb(\Y))$ に適用）：
      ボレル可測写像 $T \colon \X \to \Y$ に対し，
      \[
        \forall B \in \Bb(\Y),\quad T^{-1}(B) \in \Bb(\X).
      \]
  \end{enumerate}
\end{proposition}

%% ============================================================
%%  5. 測度に対する積分
%% ============================================================
\section{測度に対する積分}
\label{sec:sem-integration}

\begin{remark}{非負可測関数の積分}{sem-nonneg-integral}
  可測空間 $(\Omega, \mathcal{F})$ 上の確率測度 $\mu$ と
  非負可測関数 $f \colon \Omega \to [0, \infty]$ に対して，
  積分 $\int_\Omega f \, \d\mu$ は常に $[0, +\infty]$ に値を持ち
  well-defined である．
  特に，$f$ が有界であるか $\Omega$ がコンパクトであれば積分は有限値をとる．

  一般の可測関数 $f = f^+ - f^-$（$f^+ = \max(f, 0)$，$f^- = \max(-f, 0)$）に対しては，
  $\int f^+$ と $\int f^-$ の少なくとも一方が有限でなければ
  積分は定義できない（$\infty - \infty$ が不定となるため）．
  本稿で扱うコスト関数は $c \geq 0$ であるから，この問題は生じない．
\end{remark}

%% ============================================================
%%  6. Dirac 測度と押し出し
%% ============================================================
\section{Dirac 測度と押し出し}
\label{sec:sem-dirac-pushforward}

\begin{definition}{Dirac 測度}{sem-dirac}
  距離空間 $\X$ 上の点 $x \in \X$ に対して，
  \textbf{Dirac 測度} $\delta_x \in \Mm_+^1(\X)$ を
  \[
    \delta_x(A) \defeq
    \begin{cases}
      1 & x \in A, \\
      0 & x \notin A
    \end{cases}
    \qquad (\forall\, A \in \Bb(\X))
  \]
  で定義する．
  直観的には，全質量が点 $x$ に集中した確率測度である．
\end{definition}

\begin{remark}{Dirac 測度に対する積分}{sem-dirac-integral}
  非負可測関数 $f \colon \X \to [0, \infty]$ に対して，
  \[
    \int_\X f \, \d\delta_x = f(x)
  \]
  が成り立つ．以下にこれを示す．

  集合 $A \in \Bb(\X)$ に対して，
  \textbf{指示関数} $\mathbf{1}_A \colon \X \to \{0, 1\}$ を
  $\mathbf{1}_A(y) = 1$（$y \in A$），$\mathbf{1}_A(y) = 0$（$y \notin A$）
  で定める．
  $\mathbf{1}_A$ の積分は $\int \mathbf{1}_A \, \d\delta_x = \delta_x(A) = \mathbf{1}_A(x)$ である．

  有限個の可測集合と非負係数による線形結合
  $f = \sum_k c_k \mathbf{1}_{A_k}$（\textbf{単関数}）のとき，
  積分の線形性から
  \[
    \int_\X f \, \d\delta_x
    = \sum_k c_k \, \delta_x(A_k)
    = \sum_k c_k \, \mathbf{1}_{A_k}(x)
    = f(x).
  \]
  一般の非負可測関数 $f$ に対しては，
  $f_1 \leq f_2 \leq \cdots \leq f$ かつ各点で $f_n(x) \to f(x)$ を満たす
  単関数の単調増加列 $(f_n)$ が存在することが知られている．
  \textbf{単調収束定理}（非負可測関数の単調増加列に対して
  積分と極限の交換 $\int \lim f_n = \lim \int f_n$ が成り立つ）により，
  \[
    \int_\X f \, \d\delta_x
    = \lim_{n \to \infty} \int_\X f_n \, \d\delta_x
    = \lim_{n \to \infty} f_n(x)
    = f(x).
  \]
\end{remark}

\begin{definition}{押し出し（push-forward）}{sem-pushforward}
  距離空間 $\X, \Y$ の間のボレル可測写像 $T \colon \X \to \Y$ と
  測度 $\mu \in \Mm_+^1(\X)$ に対して，
  $T$ による $\mu$ の\textbf{押し出し} $T\pushforward \mu \in \Mm_+^1(\Y)$ を
  \[
    (T\pushforward \mu)(A) \defeq \mu(T^{-1}(A))
    \qquad (\forall\, A \in \Bb(\Y))
  \]
  で定義する．
\end{definition}

%% ============================================================
%%  7. 凸結合と凸集合
%% ============================================================
\section{凸結合と凸集合}
\label{sec:sem-convex}

\begin{definition}{凸結合・凸集合}{sem-convex}
  ベクトル空間 $V$ の元 $v_1, v_2 \in V$ と $t \in [0,1]$ に対して，
  $t v_1 + (1-t) v_2$ を $v_1$ と $v_2$ の\textbf{凸結合}という．
  部分集合 $S \subseteq V$ が\textbf{凸}であるとは，
  任意の $v_1, v_2 \in S$ の凸結合が再び $S$ に属すること，すなわち
  \[
    v_1, v_2 \in S,\; t \in [0,1]
    \;\Longrightarrow\;
    t v_1 + (1-t) v_2 \in S
  \]
  が成り立つことをいう．
\end{definition}
